\documentclass[10pt,a4paper]{article}

% ---------- Packages ----------
\usepackage[utf8]{inputenc}
\usepackage[T1]{fontenc}
\usepackage[french]{babel}
\usepackage{lmodern}
\usepackage[landscape,margin=1.4cm]{geometry}
\usepackage[table]{xcolor}
\usepackage{longtable}
\usepackage{array}
\usepackage{booktabs}
\usepackage{tikz}
\usepackage{hyperref}
 \usepackage{amsmath}

% ---------- Couleurs ----------
\definecolor{bleu}{HTML}{009EE0}
\definecolor{bleufonce}{HTML}{16396B}
\definecolor{gris}{HTML}{F3F3F3}
\definecolor{vert}{HTML}{2EA44F}
\definecolor{rouge}{HTML}{D7263D}
\definecolor{orange}{HTML}{F28C00}

\hypersetup{colorlinks=true, urlcolor=bleu, linkcolor=bleufonce}

% ---------- Pastilles de position ----------
\newcommand{\pastille}[1]{\tikz[baseline=-0.55ex]\fill[#1] (0,0) circle (0.9ex);}
\newcommand{\fav}{\pastille{vert}\,Favorable}
\newcommand{\opp}{\pastille{rouge}\,Opposé}
\newcommand{\inter}{\pastille{orange}\,Intermédiaire}

% ---------- Raccourcis ----------
\newcommand{\nl}{\newline}
\newcommand{\epi}{Épistémologique}
\newcommand{\axio}{Axiologique}

% ---------- Mise en forme ----------
\renewcommand{\arraystretch}{1.3}
\setlength{\tabcolsep}{4pt}
\setlength{\parindent}{0pt}
\setlength{\emergencystretch}{3em}
\pagestyle{empty}

\newcolumntype{P}[1]{>{\raggedright\arraybackslash}p{#1}}

% Ligne de titre de groupe (un groupe = un document)
% Ligne de titre de groupe (un groupe = un document), titre cliquable
\newcommand{\groupe}[2]{%
  \rowcolor{bleufonce}
  \multicolumn{6}{@{}l@{}}{\href{#1}{\textcolor{white}{\bfseries\ \,#2}}} \\
}

\begin{document}

{\color{bleufonce}\LARGE\bfseries TD2 - Tableau d'analyse des acteurs}\par\vspace{2pt}
{\color{bleu}\large Groupe 10 -- Ferracuti Filippo - Cuenca Corentin - Daviaud Yohan - Leclere Adam - Vayva Nathan}\par \vspace{2pt}
{\color{bleu}\large 18/11/2026}
\vspace{8pt}

\textbf{Controverse :} le numérique (et l'IA) est-il un allié de la transition écologique, ou son empreinte et ses effets rebond en annulent-ils les bénéfices ?

\vspace{4pt}
\textbf{Légende :} \fav\quad \opp\quad \inter\quad (reconnaît des bénéfices mais exprime des réserves).

\vspace{8pt}

\rowcolors{2}{gris}{white}
\begin{longtable}{@{}
  P{3.9cm} P{2.7cm} P{2.7cm} P{8.6cm} P{3.4cm} P{3.0cm} @{}}

% ----- En-tête première page -----
\rowcolor{bleu}
\color{white}\bfseries Acteur &
\color{white}\bfseries Type d'acteur &
\color{white}\bfseries Position &
\color{white}\bfseries Argument principal &
\color{white}\bfseries Type d'argument &
\color{white}\bfseries Nature de l'argument \\
\endfirsthead

% ----- En-tête pages suivantes -----
\rowcolor{bleu}
\color{white}\bfseries Acteur &
\color{white}\bfseries Type d'acteur &
\color{white}\bfseries Position &
\color{white}\bfseries Argument principal &
\color{white}\bfseries Type d'argument &
\color{white}\bfseries Nature de l'argument \\
\endhead

\multicolumn{6}{r}{\scriptsize\itshape Suite page suivante\ldots} \\
\endfoot
\bottomrule
\endlastfoot

% =====================================================
% DOCUMENT 1 -- Article scientifique
% =====================================================
\groupe{}{Document 1 -- Communiqué de presse : ADEME, 
\emph{Le numérique : allié de la transition écologique ?} (06/11/2025)}

Callegari, Misganaw \& Bagherzadeh (NEOMA, Kristiania)
& Scientifique & \inter
& La numérisation offre des avantages, mais elle ne suffit pas à elle seule à atteindre la neutralité carbone : elle repose sur une infrastructure matérielle et énergivore, génère des effets rebond, élargit l'offre au lieu de la substituer et ne change pas la gouvernance des entreprises. Il faut une « double transformation » (numérique et durable).
& Scientifique \nl Environnemental & \epi \\

Alibaba Group (D.~Zhang)
& Économique & \fav
& La numérisation et les plateformes décarbonent les opérations de l'entreprise et toute sa chaîne de valeur.
& Environnemental \nl Économique & \epi \\

Forum économique mondial (WEF) et Accenture
& Économique & \fav
& Le secteur numérique est l'« influenceur » le plus puissant pour accélérer l'action vers les objectifs de l'Accord de Paris, et ses technologies pourraient réduire les émissions de 20\,\% d'ici 2050 (analyse Accenture--WEF).
& Environnemental & \epi \\

Cabinets de conseil (McKinsey, PwC, EY)
& Économique & \fav
& Leurs rapports soulignent le rôle clé des technologies numériques dans la réduction des empreintes carbone.
& Environnemental \nl Économique & \epi \\

Fournisseurs de technologies ICT (Ericsson)
& Économique & \inter
& Reconnaissent l'enjeu de la consommation énergétique du secteur et s'engagent à la neutralité carbone de leurs chaînes de valeur d'ici 2050, tout en étant jugés essentiels aux objectifs de « net zero » des pays.
& Environnemental \nl Économique & \axio \\

Décideurs politiques
& Politique & \fav
& Promeuvent et financent la numérisation rapide de tous les secteurs de l'économie.
& Économique & \axio \\

Acteurs de la blockchain et grandes marques (Nike, Adidas)
& Économique & \fav
& Élargissent leur offre par de nouveaux modèles d'affaires numériques : cryptomonnaies pour les acteurs de la blockchain, NFT pour Nike, Adidas et Under Armour.
& Économique & \axio \\

\newpage
% =====================================================
% DOCUMENT 2 -- Communiqué ADEME
% =====================================================
\groupe{}{Document 2 -- Article de presse institutionnelle : \emph{Le poids écologique du numérique : comment freiner l'emballement ?}, ADEME Infos (avis de janvier 2025)}

ADEME (étude IT4Green)
& Scientifique & \inter
& Les gains environnementaux sont réels dans les 5 cas étudiés mais modestes et fragiles, car l'effet rebond et la dépendance aux métaux peuvent les annuler.
& Scientifique \nl Environnemental & \epi \\

ADEME (E.~Fangeat)
& Scientifique & \inter
& Le numérique peut être un allié à condition d'être piloté avec sobriété, au juste niveau de technologie et au service d'un besoin réel.
& Environnemental \nl Éthique & \axio \\

Ministères de l'Écologie et du Numérique
& Politique & \inter
& Face à l'empreinte carbone croissante du numérique et aux promesses de ses technologies, ils commandent l'étude IT4Green pour en évaluer objectivement les impacts nets.
& Scientifique & \epi \\

%\newpage
% =====================================================
% DOCUMENT 3 -- ADEME Infos
% =====================================================
\groupe{}{Document 3 -- Article scientifique : B.~Callegari, B.~A.~Misganaw et M.~Bagherzadeh, \emph{Paved with good intentions: From digitalization to carbon neutrality}, Journal of Environmental Management (2024)}

ADEME (R.~Guastavi, M.~Wellhoff) -- constat
& Scientifique & \opp
& L'empreinte carbone du numérique en France est passée de 17 à 29,5\,MtCO$_2$e en deux ans, en partie grâce à une meilleure prise en compte des impacts (data centers à l'étranger) et en partie sous l'effet de la multiplication des usages ; cet emballement exponentiel menace le climat, la biodiversité et les efforts de transition.
& Environnemental & \epi \\

ADEME (R.~Guastavi, M.~Wellhoff) -- recommandations
& Scientifique & \inter
& Il faut allonger la durée de vie des appareils et prioriser les usages, sans tout supprimer, car certains outils numériques servent la transition.
& Environnemental \nl Social & \axio \\

Agence internationale de l'énergie (AIE)
& Scientifique & \opp
& Prévoit une hausse de plus de 75\,\% de la consommation électrique des data centers d'ici 2026, portée par l'essor de l'IA générative (chiffre cité par l'ADEME).
& Environnemental & \epi \\

Microsoft et Google
& Économique & \inter
& Misent sur le nucléaire pour alimenter l'IA (relance de Three Mile Island pour Microsoft, mini-réacteurs SMR pour Google) tout en admettant qu'ils auront du mal à tenir leurs engagements climatiques (émissions de Microsoft : +13\,\% depuis 2020).
& Économique \nl Environnemental & \axio \\

États (non nommés dans l'article)
& Politique & \fav
& Relancent leurs centrales à charbon, en attendant la construction de centrales nucléaires, pour répondre à la hausse de la demande électrique liée à l'essor du numérique et de l'IA.
& Économique \nl Environnemental & \axio \\

\newpage
% =====================================================
% DOCUMENT 4 -- MIT Technology Review
% =====================================================

\groupe{https://www.technologyreview.com/2025/04/10/1114912/why-the-climate-promises-of-ai-sound-a-lot-like-carbon-offsets/amp/}{Document 4 -- Article de presse : J.~Temple, \emph{Why the climate promises of AI sound a lot like carbon offsets}, MIT Technology Review (10/04/2025)}


James Temple (MIT Technology Review)
& Médiatique & \inter
& Reconnaît que l'IA pourrait réduire les émissions à terme, mais justifier aujourd'hui des data centers alimentés par des fossiles par ces gains futurs ressemble aux crédits carbone : bénéfices incertains, lointains et sans mécanisme de responsabilité.
& Environnemental \nl Économique & \axio \\

Agence internationale de l'énergie (AIE)
& Scientifique & \inter
& L'IA pourrait réduire les émissions de 1,4\,Gt en 2035, mais ce scénario n'est pas une prédiction et l'impact pourrait être marginal sans conditions favorables.
& Scientifique \nl Environnemental & \epi \\

Sam Altman (OpenAI)
& Économique & \fav
& L'IA apportera des « triomphes stupéfiants », dont régler le climat, avec une intelligence quasi illimitée et une énergie abondante ; Temple y voit une justification implicite de la forte demande électrique des data centers.
& Environnemental \nl Économique & \axio \\

Développeurs de data centers et industrie gazière
& Économique & \opp
& Leur pratique (gaz, charbon reconverti) contredit la thèse du numérique comme allié de la transition.
& Économique & \axio \\

Entreprises technologiques et électriciens engagés
& Économique & \fav
& Financent des fermes solaires, aident à relancer des centrales nucléaires (Three Mile Island) ou signent des contrats géothermiques pour alimenter leurs data centers en énergie bas carbone.
& Environnemental \nl Économique & \axio \\

\newpage
% =====================================================
% DOCUMENT 5
% =====================================================

\groupe{}{Document 5 -- Article scientifique : C.~Freitag, M.~Berners-Lee, K.~Widdicks, B.~Knowles, G.~S.~Blair et A.~Friday, \emph{The real climate and transformative impact of ICT: A critique of estimates, trends, and regulations}, Patterns (10/09/2021)}

Freitag et~al. (Lancaster University / Small World Consulting)
& Scientifique & \inter
& Estiment que les TIC représentent 1,8\,\%--2,8\,\% des émissions mondiales de GES (2,1\,\%--3,9\,\% après correction de l'erreur de troncature). Sans contrainte carbone globale, l'efficacité énergétique, les renouvelables et l'« habilitation » ne suffiront pas : des effets rebond sont probables et une régulation sectorielle contraignante est nécessaire.
& Environnemental \nl Scientifique & \epi \\

Malmodin et Lundén (Ericsson / Talia)
& Scientifique & \fav
& Estiment que les émissions des TIC se sont stabilisées (1,9\,\% du total en 2020) et pourraient être divisées par deux d'ici 2030 grâce aux renouvelables ; les auteurs relèvent toutefois des données confidentielles et un risque de conflit d'intérêts, car les deux travaillent pour des opérateurs de réseaux.
& Environnemental \nl Scientifique & \epi \\

Belkhir et Elmeligi ; Andrae et Edler (Huawei)
& Scientifique & \opp
& Projettent une hausse continue des émissions des TIC, portée par le trafic de données et le nombre d'appareils, car les gains d'efficacité s'essoufflent (fin de la loi de Moore). Leurs données sont toutefois plus anciennes (2008--2012) que celles de Malmodin.
& Environnemental \nl Scientifique & \epi \\

Masanet et al. (\emph{Science})
& Scientifique & \inter
& La consommation des data centers n'a que peu augmenté (194 à 205\,TWh entre 2010 et 2020) malgré une hausse de 550\,\% des calculs, mais l'efficacité ne compensera qu'un doublement de la demande ; au-delà, la consommation croîtra vite sans action politique.
& Scientifique \nl Environnemental & \epi \\

GeSI (\emph{SMARTer 2030}, rapport sponsorisé par des entreprises des TIC)
& Économique & \fav
& Affirme que les TIC pourraient éviter 12,08\,Gt\,CO$_2$e en 2030 dans les autres secteurs, soit bien plus que leurs propres émissions, et seraient donc « carbone négatives ». Les auteurs critiquent le manque de transparence, la minimisation des effets rebond et le conflit d'intérêts.
& Environnemental \nl Économique & \epi \\

Commission européenne (Pacte vert)
& Politique & \fav
& Fait des TIC (IA, IoT, blockchain) un pilier de la neutralité climatique et vise des data centers neutres en carbone d'ici 2030, mais sans mécanisme d'incitation ou de sanction, ni estimation des réductions attendues. Cette stratégie concilie compétitivité économique et climat, au risque de stimuler la consommation.d
& Politique \nl Économique \nl Environnemental & \axio \\

Grandes entreprises des TIC (Microsoft, Amazon, BT, Sky, Apple)
& Économique & \inter
& Annoncent des engagements de neutralité carbone, de net zéro, voire de « carbone négatif » (Microsoft seule), mais ils sont inégalement ambitieux, reposent parfois sur des compensations peu additionnelles et ne couvrent pas toujours les émissions de la chaîne d'approvisionnement.
& Environnemental \nl Économique & \axio \\

Union internationale des télécommunications (UIT, avec GSMA, GeSI et SBTi)
& Institutionnel & \inter
& Propose depuis 2020 une norme volontaire visant $-45$\,\% d'émissions d'ici 2030 et le net zéro en 2050 pour le secteur. Les auteurs la jugent utile comme modèle, mais insuffisante sans mécanismes de conformité obligatoires.
& Environnemental \nl Politique & \axio \\


\groupe{https://www.tf1info.fr/environnement-ecologie/environnement-le-ticket-de-caisse-dematerialise-une-mesure-pas-si-ecolo-2257218.html}{Document 6 -- Article de presse : TF1 Info, \emph{Environnement : le ticket de caisse dématérialisé, une mesure pas si écolo} (TF1 Info)}

TF1 Info (rédaction)
& Médiatique & \opp
& La dématérialisation du ticket de caisse n'est pas automatiquement plus écologique que le papier et son impact réel reste inconnu faute d'études globales et actualisées.
& Environnemental & \epi \\

Alain Houpert (sénateur LR)
& Politique & \opp
& Le ticket dématérialisé a un bilan carbone « désastreux » (+2~g de $\text{CO}_2$ par rapport au papier) et constitue une mesure idéologique contre-productive sur le plan environnemental.
& Environnemental & \epi \\

Alexis Corbière (député LFI)
& Politique & \opp
& Le stockage dans les centres de données et la transmission énergétique d'un ticket dématérialisé rejettent 2~g de $\text{CO}_2$ de plus qu'un ticket imprimé.
& Environnemental & \epi \\

Frédéric Bordage / Collectif Green IT
& Scientifique & \inter
& Les chiffres de 2~g de $\text{CO}_2$ sont issus d'une étude interne dépassée et l'impact réel est inconnu faute d'intégrer les nouveaux écrans ; la meilleure pratique reste l'absence totale de ticket (papier ou numérique).
& Environnemental & \epi \\

GreenTicket (start-up)
& Économique & \fav
& L'émission annuelle de 30 milliards de tickets papier équivaut à 250 millions d'arbres abattus et 950 millions de litres d'eau consommés, ce qui justifie de réduire le papier.
& Environnemental & \epi \\

Gouvernement français
& Politique & \inter
& La fin de l'impression automatique évite le gaspillage tout en laissant le choix au consommateur, mais il faut veiller à éviter un transfert d'impact du papier vers le numérique et encourager l'écoconception.
& Environnemental & Axiologique \\



\end{longtable}

\end{document}